\documentclass[11pt]{article}
\usepackage[margin=1in]{geometry}
\usepackage{amsmath,amssymb,amsthm,hyperref,enumitem}
\hypersetup{colorlinks=true,linkcolor=blue,urlcolor=blue}
\newtheorem{theorem}{Theorem}
\newtheorem{corollary}{Corollary}
\newcommand{\R}{\mathbb R}
\title{An obstruction to positive Taylor cancellation\\in sums of shifted roots}
\author{Principia Math In Progress}
\date{22 September 2026}
\begin{document}
\maketitle
\noindent\textbf{MathDB 374978. Full disproof of Conjecture 5.1.}

\section*{Task 1: the source conjecture}
Korsky [1, Conjecture 5.1, p.~8] proposes that for every \(d\geq2\)
and \(k\geq1\), fixed positive integers \(A_i\) and integers \(u_i,v_i\)
can be chosen so that, along infinitely many integers \(M\to+\infty\),
\[
 S(M):=\sum_{i=1}^k A_i\bigl((M+u_i)^d+v_i\bigr)^{1/d}
 =L(M)+\lambda M^{-(dk-1)}+O(M^{-dk}),                    \tag{1}
\]
where \(L\) is an integer-valued linear polynomial and
\(\lambda\ne0\). We disprove this for every \(3\leq k\leq d\).
In particular, \(d=k=3\) is an admissible parameter pair for which
no permitted choice of coefficients and shifts exists.

\begin{theorem}
Fix \(d\geq2\), positive real weights \(A_i\), and real \(u_i,v_i\),
with at most \(d\) distinct values among the \(u_i\). Define \(S\)
as above for sufficiently large positive \(M\), and put
\[
 L_0(M)=\Bigl(\sum_i A_i\Bigr)M+\sum_i A_i u_i.
\]
If \(S(M)-L(M)=o(M^{1-2d})\) along an unbounded positive sequence,
where \(L\) is linear, then \(L=L_0\) and every \(v_i=0\).
More precisely, if the coefficients of
\(M^{1-d},M^{-d},\ldots,M^{2-2d}\) all vanish, then either
every \(v_i=0\), or
\[
 S(M)-L_0(M)=
 -\frac{d-1}{2d^2}\Bigl(\sum_i A_i v_i^2\Bigr)M^{1-2d}
      +O(M^{-2d}),                                      \tag{2}
\]
with a strictly negative leading coefficient.
\end{theorem}

\begin{proof}
For fixed real \(u,v\), the binomial expansion at zero gives
\begin{align}
 \bigl((M+u)^d+v\bigr)^{1/d}
 &=M+u+\frac vd(M+u)^{1-d}\notag\\
 &\quad-\frac{d-1}{2d^2}v^2(M+u)^{1-2d}
       +O(M^{1-3d}).                                    \tag{3}
\end{align}
The radicand and \(M+u\) are positive for large \(M\);
the real branch and the Taylor remainder are therefore valid.
The constants may depend on the fixed parameters, which suffices
even for the weaker hypothesis of the theorem.

In particular \(S-L_0=O(M^{1-d})\). If \(S-L=o(1)\) on an
unbounded sequence, the linear polynomial \(L-L_0\) tends to zero
there and hence is identically zero.

Set \(\mu_j=\sum_i A_i v_i u_i^j\). In (3) the quadratic term in
\(v\) starts only at \(M^{1-2d}\). Before that order, the coefficient
of \(M^{1-d-j}\), for \(0\leq j\leq d-1\), is precisely
\[
 \frac1d\binom{1-d}{j}\mu_j.                             \tag{4}
\]
All displayed binomial factors are nonzero, since
\(\binom{1-d}{j}=(-1)^j\binom{d+j-2}{j}\).
The hypothesis \(S-L_0=o(M^{1-2d})\) forces all these coefficients
to vanish. This remains true on a subsequence: a first nonzero
coefficient at an earlier order, divided by its power of \(M\),
would tend to a nonzero constant there. Thus
\[
 \mu_0=\mu_1=\cdots=\mu_{d-1}=0.                         \tag{5}
\]

Let \(s_1,\ldots,s_t\) be the distinct shifts, where \(t\leq d\), and
put \(c_\ell=\sum_{i:u_i=s_\ell}A_i v_i\).
The first \(t\) equations in (5) are
\[
 \sum_{\ell=1}^t c_\ell s_\ell^j=0
       \quad(0\leq j<t).
\]
The Vandermonde determinant is
\(\prod_{\ell<h}(s_h-s_\ell)\ne0\), so every \(c_\ell=0\).
This argument explicitly includes repeated shifts.
It follows that the entire first-order contribution in (3)
cancels identically:
\[
 \sum_i A_i v_i(M+u_i)^{1-d}
       =\sum_{\ell=1}^t c_\ell(M+s_\ell)^{1-d}=0.         \tag{6}
\]
Summing the remaining terms of (3), and using
\((M+u_i)^{1-2d}=M^{1-2d}(1+O(M^{-1}))\), gives (2).
Its coefficient is zero exactly when all \(v_i=0\), because
every weight is positive. Otherwise (2) contradicts
\(o(M^{1-2d})\). This proves both assertions.
\end{proof}

\begin{corollary}
Conjecture 5.1 is false for every \(3\leq k\leq d\), even if its
integer parameters are relaxed to real parameters and positive
real weights.
\end{corollary}
\begin{proof}
There are at most \(k\leq d\) distinct shifts, and
\(dk-1>2d-1\). Thus (1) implies the theorem's little-oh
hypothesis. It forces all \(v_i=0\), and hence \(S=L=L_0\)
identically. Multiplying (1) by \(M^{dk-1}\) then gives
\(0=\lambda+O(M^{-1})\), contrary to \(\lambda\ne0\).
\end{proof}

For the concrete case \(d=k=3\), the four relevant coefficients are
\[
 [M^{-2}]\,(S-L_0)=\frac{\mu_0}{3},\quad
 [M^{-3}]\,(S-L_0)=-\frac{2\mu_1}{3},\quad
 [M^{-4}]\,(S-L_0)=\mu_2,
\]
\[
 [M^{-5}]\,(S-L_0)
       =-\frac43\mu_3-\frac19\sum_i A_i v_i^2.
\]
The first three vanish only if the grouped weights \(c_\ell\)
vanish, which also makes \(\mu_3=0\). The fourth then forces
every \(v_i=0\). Already these four coefficients preclude the
proposed first nonzero term at \(M^{-8}\).

The source uses \(M\to+\infty\), as its radicands grow.
Even if arbitrary unbounded integer subsequences were permitted,
a negative subsequence gives the same obstruction after \(T=-M\):
for even \(d\) replace \(u_i\) by \(-u_i\); for odd \(d\) negate
the whole sum and replace \((u_i,v_i)\) by \((-u_i,-v_i)\).
The transformed weights remain positive and the leading coefficient
in (1) remains nonzero.

\section*{Sharpness and the separate approximation question}
The barrier order \(2d-1\) is attained with two terms:
\[
 (M^d+1)^{1/d}+(M^d-1)^{1/d}
 =2M-\frac{d-1}{d^2}M^{1-2d}+O(M^{1-4d}).               \tag{7}
\]
Thus positivity does not rule out all cancellation.
The source's square-root examples with three or four terms have
three distinct shifts, exceeding \(d=2\), and do not contradict
the theorem.

The obstruction concerns the precise fixed-shift Taylor family in
Conjecture 5.1. It does not disprove the separate bound for
\(g_{k,d}(N)\) using arbitrary integer radicands, or the
uniform approximation conjecture 1.2. The construction in (1)
would imply such an integer-target bound; the converse is not
established by the source's use of the word \emph{equivalently}.
Our result is a structural disproof of the explicit cancellation
conjecture, not credit for that wording issue.

\section*{Tasks 2--5: four proposed mathematical follow-ups}
\begin{enumerate}[label=\textbf{Task \arabic*:},start=2,leftmargin=*]
\item Determine the least number of distinct shifts permitting
nontrivial positive cancellation beyond order \(M^{1-2d}\).
The theorem proves that at least \(d+1\) are necessary.
\item Settle the square-root Taylor-cancellation assertion for
every \(k\geq5\), extending the source's \(k=2,3,4\) constructions.
\item Classify attainable leading orders when signed nonzero
rational weights are allowed in place of positive weights.
\item Resolve the separate unrestricted integer-target bound
\(g_{3,3}(N)\ll N^{-8/3}\), using radicands beyond the prohibited
three-term fixed-shift family.
\end{enumerate}
These are proposed directions and earn no additional solution credit.

\section*{Verification and source}
Four exact verification groups check the radical expansion against
its defining power equation, Vandermonde determinants with arbitrary
signs of the shifts, the cubic obstruction, and the sharp two-term
family together with the three square-root examples in [1].
The proof for arbitrary parameters is analytic and algebraic;
finite checks only supplement it. Review is by the authoring agent
only, without independent peer review or a proof assistant.

\medskip\noindent
[1] Samuel Korsky, \emph{Inhomogeneous Approximation by Sums of Roots},
\href{https://arxiv.org/abs/2605.27233}{arXiv:2605.27233v1},
26 May 2026. All ten pages were read. The target is Conjecture 5.1
on p.~8, restated as a coefficient system on p.~10.
The source's square-root constructions are attributed, not claimed
as new. Dated live MathDB checks and bounded public-claim and
publication searches are recorded separately; no earlier full
target solution or later target version was identified in them.
\end{document}
